\documentclass{szb-solution}

\area{Függvények}
\title{Folytonosság és határérték}
\subject{Matematika G1}
\subjectCode{BMETE94BG01}
\date{Utoljára frissítve: \today}
\docno{6}

\begin{document}
\maketitle

\subsection{Határérték a Heine- és a Cauchy-definíció alapján}

Tekintsük az
$$
  f(x)=\frac{3x+1}{5x+4}
$$
függvényt. Igazoljuk, hogy
$$
  \lim_{x\to2}f(x)=\frac12
  \text.
$$

\paragraph{Heine-definíció.}
Legyen $(x_n)$ tetszőleges olyan sorozat, amelyre $x_n\to2$, és
$x_n\neq2$. Mivel $5x_n+4\to14\neq0$, a nevező egy indextől kezdve nem
nulla, és a sorozatok határértékére vonatkozó műveleti szabályok szerint
\begin{align*}
  f(x_n)&=\frac{3x_n+1}{5x_n+4}\\
  &\longrightarrow\frac{3\cdot2+1}{5\cdot2+4}
  =\frac7{14}=\frac12
  \text.
\end{align*}
Ez minden megfelelő $(x_n)$ sorozatra igaz, tehát a Heine-definíció alapján
a határérték valóban $1/2$.

\paragraph{Cauchy-definíció.}
Legyen $\varepsilon>0$. Ha $|x-2|<1$, akkor $1<x<3$, így
$5x+4>9$. Ekkor
\begin{align*}
  \left|f(x)-\frac12\right|
  &=\left|\frac{2(3x+1)-(5x+4)}{2(5x+4)}\right|\\
  &=\frac{|x-2|}{2|5x+4|}
  <\frac{|x-2|}{18}
  \text.
\end{align*}
Válasszuk
$$
  \delta=\min\{1,18\varepsilon\}
  \text.
$$
Ha $0<|x-2|<\delta$, akkor a fenti becslés szerint
$$
  \left|f(x)-\frac12\right|
  <\frac{|x-2|}{18}<\varepsilon
  \text.
$$
Ez pontosan a Cauchy-féle definíció, tehát
$$
  \boxed{\displaystyle\lim_{x\to2}\frac{3x+1}{5x+4}=\frac12}
  \text.
$$

\subsection{Határértékek}

\begin{enumerate}[a)]
  \item Közös nevezőre hozva
        \begin{align*}
          &\frac{x^2}{2x+1}+\frac{x^3+4x^2-2}{1-2x^2}\\
          &\qquad=
          \frac{9x^3+5x^2-4x-2}{-4x^3-2x^2+2x+1}
          \text.
        \end{align*}
        A számláló és a nevező azonos fokú, ezért a vezető együtthatók
        hányadosa adja a határértéket:
        $$
          \boxed{\displaystyle
          \lim_{x\to-\infty}\left(
          \frac{x^2}{2x+1}+\frac{x^3+4x^2-2}{1-2x^2}
          \right)=-\frac94}
          \text.
        $$

  \item A két egyoldali határértéket külön kell vizsgálni. Ha $x\to0^+$,
        akkor $\sqrt2/x\to+\infty$, míg $x\to0^-$ esetén
        $\sqrt2/x\to-\infty$. Így
        $$
          \lim_{x\to0^+}\frac3{2^{\sqrt2/x}+1}=0,
          \qquad
          \lim_{x\to0^-}\frac3{2^{\sqrt2/x}+1}=3
          \text.
        $$
        Mivel a két egyoldali határérték különbözik,
        $$
          \boxed{\displaystyle
          \lim_{x\to0}\frac3{2^{\sqrt2/x}+1}\text{ nem létezik}}
          \text.
        $$

  \item Konjugálttal bővítve, majd $x+1$-gyel egyszerűsítve:
        \begin{align*}
          \frac{x+1}{\sqrt{6x^2+3}+3x}
          &=\frac{(x+1)(\sqrt{6x^2+3}-3x)}{3(1-x^2)}\\
          &=\frac{\sqrt{6x^2+3}-3x}{3(1-x)}
          \qquad(x\neq-1)
          \text.
        \end{align*}
        Ezért
        $$
          \boxed{\displaystyle
          \lim_{x\to-1}\frac{x+1}{\sqrt{6x^2+3}+3x}=1}
          \text.
        $$

  \item A kitevőben
        $$
          \frac{1-\sqrt x}{1-x}
          =\frac1{1+\sqrt x}\longrightarrow\frac12
          \text,
        $$
        az alap pedig $(1+x)/(2+x)\to2/3>0$. A hatványfüggvény
        folytonossága miatt
        $$
          \boxed{\displaystyle
          \lim_{x\to1}\left(\frac{1+x}{2+x}\right)^{
          \frac{1-\sqrt x}{1-x}}
          =\left(\frac23\right)^{1/2}=\sqrt{\frac23}}
          \text.
        $$

  \item Legyen $y=\sin x$. Ekkor $y\to1/2$, és
        $$
          2y^2+y-1=(2y-1)(y+1),
          \qquad
          2y^2-3y+1=(2y-1)(y-1)
          \text.
        $$
        Az eltűnő közös tényező egyszerűsítése után
        $$
          \boxed{\displaystyle
          \lim_{x\to\pi/6}
          \frac{2\sin^2x+\sin x-1}{2\sin^2x-3\sin x+1}
          =\frac{1/2+1}{1/2-1}=-3}
          \text.
        $$

  \item Használjuk az $1-\sin x=\cos^2x/(1+\sin x)$ azonosságot:
        \begin{align*}
          \frac{\cos x-\sin x+1}{\cos x+\sin x-1}
          &=\frac{\cos x+(1-\sin x)}{\cos x-(1-\sin x)}\\
          &=\frac{1+\dfrac{\cos x}{1+\sin x}}
                  {1-\dfrac{\cos x}{1+\sin x}}
          \longrightarrow1
          \text.
        \end{align*}
        Tehát
        $$
          \boxed{\displaystyle
          \lim_{x\to\pi/2}\frac{\cos x-\sin x+1}
          {\cos x+\sin x-1}=1}
          \text.
        $$

  \item Az alap-határértéket felhasználva
        $$
          \boxed{\displaystyle
          \lim_{x\to0}\frac{\sin5x}{x}
          =5\lim_{x\to0}\frac{\sin5x}{5x}=5}
          \text.
        $$

  \item Átalakítással minden tényező ismert határértékű:
        \begin{align*}
          \frac{\tan x-\sin x}{x^3}
          &=\frac{\sin x(1-\cos x)}{x^3\cos x}\\
          &=\frac{\sin x}{x}\cdot
            \frac{1-\cos x}{x^2}\cdot\frac1{\cos x}
          \longrightarrow1\cdot\frac12\cdot1
          \text.
        \end{align*}
        Így a határérték
        $$
          \boxed{\frac12}
          \text.
        $$

  \item
        \begin{align*}
          \frac{\sin8x}{\tan5x}
          &=\frac{\sin8x}{8x}\cdot
            \frac{5x}{\sin5x}\cdot
            \frac85\cos5x
          \longrightarrow1\cdot1\cdot\frac85\cdot1
          \text.
        \end{align*}
        Tehát
        $$
          \boxed{\displaystyle\lim_{x\to0}\frac{\sin8x}{\tan5x}=\frac85}
          \text.
        $$

  \item Legyen $t=\pi/2-x$. Ekkor $t\to0$, és
        $$
          \left(\frac\pi2-x\right)\tan x
          =t\tan\left(\frac\pi2-t\right)
          =\frac{t\cos t}{\sin t}
          \longrightarrow1
          \text.
        $$
        Ezért a határérték $\boxed{1}$.

  \item A $t=\sin x$ helyettesítéssel $t\to0$, ezért
        $$
          \boxed{\displaystyle
          \lim_{x\to0}\frac{1-\cos(\sin x)}{\sin^2x}
          =\lim_{t\to0}\frac{1-\cos t}{t^2}=\frac12}
          \text.
        $$
\end{enumerate}

\subsection{Racionális függvény folytonossága}

Legyen
$$
  f(x)=\frac{x^2+3x-10}{x^2-3x+2}
  =\frac{(x-2)(x+5)}{(x-1)(x-2)}
  \text.
$$
Az eredeti függvény értelmezési tartománya
$\mathbb R\setminus\{1,2\}$. Ezen belül racionális függvényként mindenütt
folytonos. Ha $x\neq1,2$, akkor
$$
  f(x)=\frac{x+5}{x-1}
  \text.
$$

Az $x=1$ pontban
$$
  \lim_{x\to1^-}f(x)=-\infty,
  \qquad
  \lim_{x\to1^+}f(x)=+\infty
  \text,
$$
tehát itt nem megszüntethető szakadási hely, hanem függőleges aszimptota van.
Az $x=2$ pontban viszont
$$
  \lim_{x\to2}f(x)=\lim_{x\to2}\frac{x+5}{x-1}=7
  \text.
$$
Ez megszüntethető szakadás: ha $f(2)=7$-et definiálnánk, a függvény ott is
folytonossá válna. Összefoglalva:
$$
  \boxed{f\text{ folytonos }\mathbb R\setminus\{1,2\}\text{-n};
  \ x=2\text{ megszüntethető, }x=1\text{ nem megszüntethető szakadás}.}
$$

\subsection{Egy oszcilláló függvény a nullában}

Az $f(x)=x\sin(1/x)$ képlet $x=0$-ban önmagában nem értelmezett, ezért az
eredeti függvény ott nem lehet folytonos. Vizsgáljuk azonban a természetes
kiterjesztést:
$$
  \widetilde f(x)=
  \begin{cases}
    x\sin(1/x),&x\neq0,\\
    0,&x=0.
  \end{cases}
$$
Mivel $|\sin(1/x)|\leq1$,
$$
  |x\sin(1/x)|\leq|x|\longrightarrow0
  \qquad(x\to0)
  \text.
$$
A rendőrelv szerint $\lim_{x\to0}x\sin(1/x)=0=\widetilde f(0)$, így
$$
  \boxed{\widetilde f\text{ folytonos a }0\text{ pontban}.}
$$

\subsection{Paraméteres, szakaszonként definiált függvény}

Legyen
$$
  f(x)=
  \begin{cases}
    x,&|x|<1,\\
    x^2+ax+b,&|x|\geq1.
  \end{cases}
$$
A két ág külön-külön folytonos, ezért csak $x=-1$ és $x=1$ illesztését
kell megvizsgálni. Az $x=1$ pontban
$$
  1+a+b=1,
  \qquad\text{vagyis}\qquad a+b=0
  \text.
$$
Az $x=-1$ pontban
$$
  1-a+b=-1,
  \qquad\text{vagyis}\qquad -a+b=-2
  \text.
$$
A két egyenlet megoldása
$$
  \boxed{a=1,\qquad b=-1}
  \text.
$$
Ezekkel az értékekkel a függvény mindkét csatlakozási pontban, tehát az
egész valós számegyenesen folytonos.

\subsection{Összetettebb határértékek}

\begin{enumerate}[a)]
  \item A szorzatalakhoz használjuk a
        $\lim_{u\to0}\ln(1+u)/u=1$ határértéket:
        \begin{align*}
          \frac{\ln(\cos2x)}{x^2}
          &=\frac{\ln(\cos2x)}{\cos2x-1}
            \cdot\frac{\cos2x-1}{x^2}\\
          &\longrightarrow1\cdot(-2)=-2
          \text,
        \end{align*}
        hiszen $1-\cos2x=2\sin^2x$ és $\sin x/x\to1$. Tehát
        $$
          \boxed{\displaystyle\lim_{x\to0}\frac{\ln(\cos2x)}{x^2}=-2}
          \text.
        $$

  \item A kifejezés $1^{\infty}$ típusú. Legyen
        $$
          y=(\tan x)^{\tan2x}
          \text,
        $$
        ami $x$-nek $\pi/4$-hez elég közeli, de attól különböző értékein
        pozitív. Logaritmust véve
        $$
          \ln y=\tan2x\,\ln(\tan x)
          \text.
        $$
        Tegyük $t=\tan x$-et; ekkor $t\to1$, továbbá
        $$
          \tan2x=\frac{2t}{1-t^2}
          =-\frac{2t}{(t-1)(t+1)}
          \text.
        $$
        Így
        \begin{align*}
          \lim_{x\to\pi/4}\ln y
          &=-\lim_{t\to1}
            \frac{2t}{t+1}\cdot\frac{\ln t}{t-1}\\
          &=-1\cdot1=-1
          \text.
        \end{align*}
        Az exponenciális függvény folytonosságából
        $$
          \boxed{\displaystyle
          \lim_{x\to\pi/4}(\tan x)^{\tan2x}=e^{-1}}
          \text.
        $$
\end{enumerate}

\end{document}
